% GENERATED FILE. Edit canonical lessons, then run npm run generate:books.
% canonical-source-hash: fnv1a64:1820e1466a205c1d
% canonical-lessons: FR-C521-lecteur, FR-C521-narrateur, FR-C521-heros, FR-C521-dedicace, FR-C521-dedicacer, FR-R521-revision-514-521

\chapter{Auteurs et lecteurs}
\label{ch:fr-auteurs-et-lecteurs}
\hlchaptermodality{521}

\begin{chapteropening}
\textbf{By the end of this chapter:} \emph{I can talk about readers, heroes and signed books.}

Complete the last lesson of chapter 521: I can talk about readers, heroes and signed books.
\end{chapteropening}

\section[le lecteur]{le lecteur — reader; player (for discs)}

\label{lesson:FR-C521-lecteur}



\begin{quote}
Before the new one: say the French for paperback, then the French for to devour.
\end{quote}

\subsection*{You'll want to know: le lecteur}
\textbf{le lecteur} — \textquotedblleft{}reader; player (for discs)\textquotedblright{}.

\begin{quote}
Ce journal a beaucoup de lecteurs jeunes.
\end{quote}

\begin{quote}
\emph{This newspaper has a lot of young readers.}
\end{quote}

A woman is \emph{la lectrice}. \emph{Un lecteur de DVD} — a DVD player.

\begin{cousinweb}
\emph{le lecteur} is a learned borrowing from Latin \emph{lēctor} \textquotedblleft{}a reader\textquotedblright{} (EtymDB), from \emph{legere} \textquotedblleft{}to gather, to read\textquotedblright{} (Lewis \& Short).
\end{cousinweb}

\begin{grammarlens}[title={What you've built}]
One more word for telling, explaining and describing at B1.
\end{grammarlens}

\subsection*{Your turn}
\begin{itemize}
\raggedright
  \item \emph{Say it:} \emph{le lecteur}
  \item \emph{Say it:} \emph{le lecteur}, once more
  \item \emph{Say it:} a sentence of your own with \emph{le lecteur}
  \item \emph{Recall:} say the French for paperback, then the French for to devour, then say \emph{le lecteur} again
\end{itemize}

\begin{tcolorbox}[breakable,colback=teal!4,colframe=teal!35!black,title={Before you move on}]
What does le lecteur mean? (\textquotedblleft{}Reader; player (for discs)\textquotedblright{}.) Say it once more.
\end{tcolorbox}
\section[le narrateur]{le narrateur — narrator}

\label{lesson:FR-C521-narrateur}



\begin{quote}
Before the new one: say the French for to devour, then the French for reader.
\end{quote}

\subsection*{You'll want to know: le narrateur}
\textbf{le narrateur} — \textquotedblleft{}narrator\textquotedblright{}.

\begin{quote}
Dans ce roman, le narrateur est un enfant.
\end{quote}

\begin{quote}
\emph{In this novel the narrator is a child.}
\end{quote}

A woman is \emph{la narratrice}.

\begin{cousinweb}
\emph{le narrateur} is a learned borrowing from Latin \emph{narrātor} \textquotedblleft{}a storyteller\textquotedblright{} (EtymDB), from \emph{narrāre} \textquotedblleft{}to tell\textquotedblright{} (Lewis \& Short).
\end{cousinweb}

\begin{grammarlens}[title={What you've built}]
One more word for telling, explaining and describing at B1.
\end{grammarlens}

\subsection*{Your turn}
\begin{itemize}
\raggedright
  \item \emph{Say it:} \emph{le narrateur}
  \item \emph{Say it:} \emph{le narrateur}, once more
  \item \emph{Say it:} a sentence of your own with \emph{le narrateur}
  \item \emph{Recall:} say the French for to devour, then the French for reader, then say \emph{le narrateur} again
\end{itemize}

\begin{tcolorbox}[breakable,colback=teal!4,colframe=teal!35!black,title={Before you move on}]
What does le narrateur mean? (\textquotedblleft{}Narrator\textquotedblright{}.) Say it once more.
\end{tcolorbox}
\section[le héros]{le héros — hero; main character}

\label{lesson:FR-C521-heros}



\begin{quote}
Before the new one: say the French for reader, then the French for narrator.
\end{quote}

\subsection*{You'll want to know: le héros}
\textbf{le héros} — \textquotedblleft{}hero; main character\textquotedblright{}.

\begin{quote}
Le héros du livre est un vieux marin.
\end{quote}

\begin{quote}
\emph{The hero of the book is an old sailor.}
\end{quote}

The \emph{h} is aspirated: \emph{le héros}, \emph{les héros} with no liaison.

\begin{cousinweb}
\emph{le héros} comes from Latin \emph{hērōs}, a borrowing from Greek \emph{h\'{\={e}}rōs} \textquotedblleft{}a demigod, a hero\textquotedblright{} (EtymDB).
\end{cousinweb}

\begin{grammarlens}[title={What you've built}]
One more word for telling, explaining and describing at B1.
\end{grammarlens}

\subsection*{Your turn}
\begin{itemize}
\raggedright
  \item \emph{Say it:} \emph{le héros}
  \item \emph{Say it:} \emph{le héros}, once more
  \item \emph{Say it:} a sentence of your own with \emph{le héros}
  \item \emph{Recall:} say the French for reader, then the French for narrator, then say \emph{le héros} again
\end{itemize}

\begin{tcolorbox}[breakable,colback=teal!4,colframe=teal!35!black,title={Before you move on}]
What does le héros mean? (\textquotedblleft{}Hero; main character\textquotedblright{}.) Say it once more.
\end{tcolorbox}
\section[la dédicace]{la dédicace — signed message, inscription (in a book); dedication}

\label{lesson:FR-C521-dedicace}



\begin{quote}
Before the new one: say the French for narrator, then the French for hero.
\end{quote}

\subsection*{You'll want to know: la dédicace}
\textbf{la dédicace} — \textquotedblleft{}signed message, inscription (in a book); dedication\textquotedblright{}.

\begin{quote}
L'auteur a écrit une dédicace dans mon livre.
\end{quote}

\begin{quote}
\emph{The author wrote an inscription in my book.}
\end{quote}

\emph{Une séance de dédicaces} — a book signing.

\begin{cousinweb}
\emph{la dédicace} is a learned borrowing from Latin \emph{dēdicātiō} \textquotedblleft{}a consecration\textquotedblright{}, from \emph{dēdicāre} \textquotedblleft{}to dedicate\textquotedblright{} (Lewis \& Short).
\end{cousinweb}

\begin{grammarlens}[title={What you've built}]
One more word for telling, explaining and describing at B1.
\end{grammarlens}

\subsection*{Your turn}
\begin{itemize}
\raggedright
  \item \emph{Say it:} \emph{la dédicace}
  \item \emph{Say it:} \emph{la dédicace}, once more
  \item \emph{Say it:} a sentence of your own with \emph{la dédicace}
  \item \emph{Recall:} say the French for narrator, then the French for hero, then say \emph{la dédicace} again
\end{itemize}

\begin{tcolorbox}[breakable,colback=teal!4,colframe=teal!35!black,title={Before you move on}]
What does la dédicace mean? (\textquotedblleft{}Signed message, inscription (in a book); dedication\textquotedblright{}.) Say it once more.
\end{tcolorbox}
\section[dédicacer]{dédicacer — to sign (a book for someone)}

\label{lesson:FR-C521-dedicacer}



\begin{quote}
Before the new one: say the French for hero, then the French for signed message.
\end{quote}

\subsection*{You'll want to know: dédicacer}
\textbf{dédicacer} — \textquotedblleft{}to sign (a book for someone)\textquotedblright{}.

\begin{quote}
L'écrivain dédicace ses livres à la librairie samedi.
\end{quote}

\begin{quote}
\emph{The writer is signing his books at the bookshop on Saturday.}
\end{quote}

Spelling: \emph{nous dédicaçons}.

\begin{cousinweb}
\emph{dédicacer} is built on \emph{dédicace}, a learned borrowing from Latin \emph{dēdicātiō}, from \emph{dēdicāre} \textquotedblleft{}to dedicate\textquotedblright{} (Lewis \& Short).
\end{cousinweb}

\begin{grammarlens}[title={What you've built}]
One more word for telling, explaining and describing at B1.
\end{grammarlens}

\subsection*{Your turn}
\begin{itemize}
\raggedright
  \item \emph{Say it:} \emph{dédicacer}
  \item \emph{Say it:} \emph{dédicacer}, once more
  \item \emph{Say it:} a sentence of your own with \emph{dédicacer}
  \item \emph{Recall:} say the French for hero, then the French for signed message, then say \emph{dédicacer} again
\end{itemize}

\begin{tcolorbox}[breakable,colback=teal!4,colframe=teal!35!black,title={Before you move on}]
What does dédicacer mean? (\textquotedblleft{}To sign (a book for someone)\textquotedblright{}.) Say it once more.
\end{tcolorbox}
\section[(dialogue)]{Retour sur les chapitres 514–521}

\label{lesson:FR-R521-revision-514-521}



\begin{quote}
Say the French for (theatre) company and wings.
\end{quote}

\subsection*{Your turn}
Say these aloud:

\begin{itemize}
\raggedright
  \item (theatre) company, wings, set, costume, interval
  \item stage fright, line, to perform, actor, applause
  \item paintbrush, canvas, watercolour, easel, to sketch
  \item sculpture, to sculpt, clay, marble, sculptor
  \item exhibition, private view, gallery, work, to put in a frame
  \item novelist, plot, ending, twist, gripping
  \item bookseller, section, to leaf through, paperback, to devour
  \item reader, narrator, hero, signed message, to sign
\end{itemize}

\begin{tcolorbox}[breakable,colback=teal!4,colframe=teal!35!black,title={Before you move on}]
(theatre) company? (\textbf{la troupe}.) Wings? (\textbf{les coulisses}.) Set? (\textbf{le décor}.) Costume? (\textbf{le costume}.) Interval? (\textbf{l'entracte}.) Stage fright? (\textbf{le trac}.) Line? (\textbf{la réplique}.) To perform? (\textbf{interpréter}.) Actor? (\textbf{le comédien}.) Applause? (\textbf{les applaudissements}.) Paintbrush? (\textbf{le pinceau}.) Canvas? (\textbf{la toile}.) Watercolour? (\textbf{l'aquarelle}.) Easel? (\textbf{le chevalet}.) To sketch? (\textbf{esquisser}.) Sculpture? (\textbf{la sculpture}.) To sculpt? (\textbf{sculpter}.) Clay? (\textbf{l'argile}.) Marble? (\textbf{le marbre}.) Sculptor? (\textbf{le sculpteur}.) Exhibition? (\textbf{l'exposition}.) Private view? (\textbf{le vernissage}.) Gallery? (\textbf{la galerie}.) Work? (\textbf{l'œuvre}.) To put in a frame? (\textbf{encadrer}.) Novelist? (\textbf{le romancier}.) Plot? (\textbf{l'intrigue}.) Ending? (\textbf{le dénouement}.) Twist? (\textbf{le rebondissement}.) Gripping? (\textbf{captivant}.) Bookseller? (\textbf{le libraire}.) Section? (\textbf{le rayon}.) To leaf through? (\textbf{feuilleter}.) Paperback? (\textbf{le livre de poche}.) To devour? (\textbf{dévorer}.) Reader? (\textbf{le lecteur}.) Narrator? (\textbf{le narrateur}.) Hero? (\textbf{le héros}.) Signed message? (\textbf{la dédicace}.) To sign? (\textbf{dédicacer}.)
\end{tcolorbox}
